\chapter{Galois Coarsening}
\label{ch:coarsening}

$\Fz$ eigenvalues live in $\GEI$ with Galois group
$\mathrm{Gal}(\Q(\sqrt5,\sqrt{-3})/\Q) \cong (\Z/2)^2$.
The three nontrivial Galois elements define canonical projections
with no free parameters.

\section{Eigenvalues as \texorpdfstring{$\GEI$}{Z[φ,ζ₆]} Elements}

\begin{definition}[Channel decomposition]
\label{def:from-channels}
\lean{FUST.Coarsening.fromChannels}
Each $\Fz$ eigenvalue decomposes as
$\lambda = \alpha \cdot \afcoeff_{\mathrm{gei}} + \beta$,
where $\alpha, \beta \in \Z[\vp]$.
\end{definition}

\begin{theorem}[$\tau$-norm formula]
\label{thm:tau-norm-formula}
\lean{FUST.Coarsening.tauNormSq_formula}
$|\lambda|^2 = \beta^2 + 12\alpha^2$,
where $12 = |\afcoeff_{\mathrm{gei}}|^2$ and
$\alpha, \beta$ are the channel components.
This is the $\tau$-norm squared, not a mass.
\end{theorem}

\begin{definition}[$\Fz$ eigenvalues in $\GEI$]
\label{def:eigenGEI}
\lean{FUST.Coarsening.eigenFζ_mod1, FUST.Coarsening.eigenFζ_mod5}
For $n \equiv 1 \pmod{6}$:
$\lambda_n = \mathrm{fromChannels}(5\alpha_n, 6\beta_n)$.
For $n \equiv 5 \pmod{6}$: AF sign negated.
\end{definition}

\begin{theorem}[Eigenvalue bridge]
\label{thm:eigenvalue-bridge}
\lean{FUST.Coarsening.Fζ_eigenvalue_bridge_mod1, FUST.Coarsening.Fζ_eigenvalue_bridge_mod5}
$\Fz(z^{6k+1})(1) = \mathrm{toComplex}(\lambda_{6k+1})$ and
$\Fz(z^{6k+5})(1) = \mathrm{toComplex}(\lambda_{6k+5})$.
\end{theorem}

\section{Coarsening Projections}

The three nontrivial Galois elements define canonical projections:
\begin{align*}
  N_\tau &: \GEI \to \Z[\vp] &\quad&\text{(remove gauge $\zs$, keep scale $\vp$)},\\
  N_\sigma &: \GEI \to \Z[\zs] &\quad&\text{(remove scale $\vp$, keep gauge $\zs$)},\\
  N &: \GEI \to \Z &\quad&\text{(remove both)}.
\end{align*}
These satisfy $N = N_{\Z[\vp]} \circ N_\tau = N_{\Z[\zs]} \circ N_\sigma$.

\begin{definition}[$N_\tau$: remove gauge, keep scale]
\label{def:tau-coarsen}
\lean{FUST.Coarsening.tauCoarsen}
$N_\tau: \GEI \to \Z[\vp]$, projection via $\tau$-norm.
\end{definition}

\begin{definition}[$N_\sigma$: remove scale, keep gauge]
\label{def:sigma-coarsen}
\lean{FUST.Coarsening.sigmaCoarsen}
$N_\sigma: \GEI \to \Z[\zs]$, projection via $\sigma$-norm.
\end{definition}

\begin{definition}[$N$: full norm]
\label{def:full-norm}
\lean{FUST.Coarsening.fullNorm}
$N: \GEI \to \Z$, projection via full Galois norm.
\end{definition}

\section{Composite States}

\begin{theorem}[Composite channel formula]
\label{thm:composite-channels}
\lean{FUST.Coarsening.composite_channels}
For eigenvalues $\lambda_1 = (\alpha_1, \beta_1)$ and $\lambda_2 = (\alpha_2, \beta_2)$:
\[
  \lambda_1 \cdot \lambda_2 =
  \mathrm{fromChannels}(\alpha_1\beta_2 + \alpha_2\beta_1,\;
  \beta_1\beta_2 - 12\,\alpha_1\alpha_2).
\]
\end{theorem}

\begin{theorem}[Active $\times$ active $=$ kernel]
\label{thm:active-product-kernel}
\lean{FUST.Coarsening.active_product_kernel}
If $a \equiv 1$ or $5$ and $b \equiv 1$ or $5 \pmod{6}$,
then $a + b$ is a kernel mode.
\end{theorem}

\begin{theorem}[Emergence $\tau$-norm]
\label{thm:emergence-tau-norm}
\lean{FUST.Coarsening.emergence_tauNormSq_matter, FUST.Coarsening.emergence_tauNormSq_antimatter}
$|\delta_{\Fz}(z^{6j+3}, z^{6k+4})(1)|^2 = |\lambda_{6(j+k+1)+1}|^2$ (matter),
$|\delta_{\Fz}(z^{6j+2}, z^{6k+3})(1)|^2 = |\lambda_{6(j+k)+5}|^2$ (antimatter).
\end{theorem}

\section{Poincar\'{e} Casimir Decomposition}

The $\tau$-norm squared $|\lambda|^2$ decomposes into the Poincar\'{e}
Casimir invariant $m^2 = P^\mu P_\mu$ plus spatial contributions.

\begin{theorem}[$\phiS$ spatial decomposition]
\label{thm:phiS-decomp}
\lean{FUST.Coarsening.phiS_decomp}
\[
  \phiS(n) = 10\,\sigma_5(n) + 5\,\sigma_3(n) + (6 - 2\vp)\,\sigma_2(n),
\]
where $\sigma_5, \sigma_3, \sigma_2$ are the three spatial $\Dz$ components.
\end{theorem}

\begin{theorem}[$\phiA = \Dz$ temporal component]
\label{thm:phiA-temporal}
\lean{FUST.Coarsening.phiA_eq_temporal}
$\phiA(n) = \Phi_{A\text{-coeff}}(n)$: the AF channel eigenvalue
equals the temporal $\Dz$ component.
\end{theorem}

\begin{theorem}[Casimir component expansion]
\label{thm:casimir-components}
\lean{FUST.Coarsening.casimir_components}
\[
  m^2(s) = \Phi_{A\text{-coeff}}(s)^2 - \sigma_5(s)^2 - \sigma_3(s)^2 - \sigma_2(s)^2.
\]
\end{theorem}

\begin{theorem}[$\tau$-norm vs Casimir]
\label{thm:tau-norm-casimir}
\lean{FUST.Coarsening.tauNormSq_casimir_relation_mod1}
For $n = 6k+1$:
\[
  |\lambda_n|^2 = 300\,m^2(n)
    + 300\bigl(\sigma_5(n)^2 + \sigma_3(n)^2 + \sigma_2(n)^2\bigr)
    + 36\,\phiS(n)^2.
\]
The coefficient $300 = 12 \times 25 = |\afcoeff|^2 \times 5^2$
is determined by the $\Fz$ structure.
The Casimir $m^2$ is the physical mass squared;
$|\lambda|^2$ includes additional spatial contributions.
\end{theorem}
